%======================================================================
\section{The density operator}\label{sec:density}
%======================================================================

\begin{objectives}
  You will describe ensembles of states and subsystems of entangled systems with density operators; recognise pure and mixed states and place qubit states in the Bloch ball; see a state as a positive, normalised linear functional on observables (the Riesz theorem for operators); and compute reduced density operators with the partial trace.
\end{objectives}

\begin{recall}
  Positive operators (Section~\ref{sec:measurement}); the trace and the Hilbert--Schmidt inner product $\Tr(A^\dagger B)$ (Sections~\ref{sec:vector-spaces} and~\ref{sec:hilbert}); the Riesz theorem (Section~\ref{sec:state-evolution}); tensor products and measurements on one part of a system (Section~\ref{sec:composite}).
\end{recall}

So far a system has been described by a state vector, and that is enough if we know exactly how it was prepared. Two situations break this assumption. In the first, the preparation is itself random: a source emits $\ket{\psi_i}$ with probability $p_i$ and we do not know which.
In the second, the system is one \emph{part} of an entangled pair: we shall see that it does not have a state vector of its own. In both cases the right tool is the density operator.

\subsection{Ensembles of states}

\begin{definition}[Ensemble and density operator]\label{def:density}
  An \term{ensemble} $\{p_i,\ket{\psi_i}\}$ is a collection of states $\ket{\psi_i}$ with probabilities $p_i\ge0$, $\sum_ip_i=1$. Its \term{density operator} is
  \[
    \rho=\sum_ip_i\ket{\psi_i}\!\bra{\psi_i}.
  \]
  If a single $p_i=1$ the state is \term{pure}, $\rho=\ket\psi\!\bra\psi$; otherwise it is \term{mixed}.
\end{definition}

Why is this the right object? Suppose the system is in $\ket{\psi_i}$ with probability $p_i$ and we measure $\{M_m\}$.
\begin{steps}
  \item \emph{Probabilities.} By total probability, $p(m)=\sum_ip_i\,p(m\mid i)=\sum_ip_i\bra{\psi_i}M_m^\dagger M_m\ket{\psi_i}=\Tr\bigl(M_m^\dagger M_m\,\rho\bigr)$, using $\bra\psi A\ket\psi=\Tr(A\ket\psi\!\bra\psi)$.
  \item \emph{After the measurement.} If the state was $\ket{\psi_i}$ it becomes $M_m\ket{\psi_i}$ (normalised), so the post-measurement ensemble has density operator $M_m\rho M_m^\dagger/p(m)$.
  \item \emph{Evolution.} Under a unitary each $\ket{\psi_i}\mapsto U\ket{\psi_i}$, so $\rho\mapsto U\rho U^\dagger$.
\end{steps}
We have shown the following translation of the postulates.

\begin{proposition}[Postulates for density operators]\label{prop:postulates-rho}
  In terms of $\rho$: (2) a unitary evolution takes $\rho\mapsto U\rho U^\dagger$; (3) a measurement $\{M_m\}$ has outcome probabilities $p(m)=\Tr(M_m^\dagger M_m\rho)$ and post-measurement states $\rho_m=M_m\rho M_m^\dagger/p(m)$; and the expected value of an observable $M$ is $\Tr(M\rho)$.
\end{proposition}

The crucial point is that \emph{the density operator contains everything that can be observed}. Different ensembles can have the same $\rho$, and then they are physically indistinguishable.

\begin{example}\label{ex:mixed-vs-superposition}
  Compare the ensemble $\bigl\{\tfrac12,\ket0;\ \tfrac12,\ket1\bigr\}$, the ensemble $\bigl\{\tfrac12,\ket+;\ \tfrac12,\ket-\bigr\}$, and the pure state $\ket+$.
\end{example}
\begin{solution}
  \begin{steps}
    \item \emph{First ensemble.} $\rho=\tfrac12\ket0\!\bra0+\tfrac12\ket1\!\bra1=\tfrac12\Id$.
    \item \emph{Second ensemble.} $\rho=\tfrac12\ket+\!\bra++\tfrac12\ket-\!\bra-=\tfrac12\Id$ as well, by the completeness relation. So the two ensembles are physically identical: no measurement distinguishes them, in spite of the fact that their members look very different.
    \item \emph{The pure state.} $\ket+\!\bra+=\tfrac12\begin{pmatrix}1&1\\1&1\end{pmatrix}\neq\tfrac12\Id$. It can be told apart from the mixtures: measuring $X$ gives $+1$ with certainty on $\ket+$, but with probability $\tfrac12$ on $\tfrac12\Id$.
  \end{steps}
  The \emph{superposition} $\ket+$ is not the same as a \emph{coin toss} between $\ket0$ and $\ket1$.
\end{solution}

\begin{example}[A three-state ensemble]\label{ex:three-ensemble}
  A source emits $\ket0$ with probability $\tfrac12$, $\ket1$ with probability $\tfrac14$ and $\ket+$ with probability $\tfrac14$. Find $\rho$, and the probability that a measurement of $X$ gives $+1$.
\end{example}
\begin{solution}
  \begin{steps}
    \item \emph{The density matrix.} $\rho=\tfrac12\begin{pmatrix}1&0\\0&0\end{pmatrix}+\tfrac14\begin{pmatrix}0&0\\0&1\end{pmatrix}+\tfrac14\cdot\tfrac12\begin{pmatrix}1&1\\1&1\end{pmatrix}=\tfrac18\begin{pmatrix}5&1\\1&3\end{pmatrix}$.
    \item \emph{Trace check.} $\Tr\rho=\tfrac18(5+3)=1$.
    \item \emph{The probability.} $P(X=+1)=\Tr(\ket+\!\bra+\rho)=\bra+\rho\ket+=\tfrac12\cdot\tfrac18(5+1+1+3)=\tfrac58$.
    \item \emph{Check by total probability.} $\tfrac12\cdot\tfrac12+\tfrac14\cdot\tfrac12+\tfrac14\cdot1=\tfrac58$.
  \end{steps}
\end{solution}

\subsection{Characterising density operators}

We can now forget ensembles and characterise density operators intrinsically.

\begin{theorem}[Characterisation]\label{thm:density-char}
  An operator $\rho$ is the density operator of some ensemble if and only if (a) $\Tr\rho=1$, and (b) $\rho$ is positive (Definition~\ref{def:positive}).
\end{theorem}
\begin{proof}
  \begin{steps}
    \item \emph{Necessity.} If $\rho=\sum_ip_i\ket{\psi_i}\!\bra{\psi_i}$, then $\Tr\rho=\sum_ip_i=1$ and $\bra\varphi\rho\ket\varphi=\sum_ip_i\abs{\braket\varphi{\psi_i}}^2\ge0$.
    \item \emph{Sufficiency.} Suppose (a) and (b). By Theorem~\ref{thm:positive}, $\rho$ is Hermitian with eigenvalues $\lambda_j\ge0$: $\rho=\sum_j\lambda_j\ket j\!\bra j$ for an orthonormal basis $\ket j$.
    \item \emph{The ensemble.} $\sum_j\lambda_j=\Tr\rho=1$, so $\{\lambda_j,\ket j\}$ is an ensemble, and its density operator is $\rho$. \qedhere
  \end{steps}
\end{proof}

When is a state pure? Since $\Tr\rho^2=\sum_j\lambda_j^2\le\bigl(\sum_j\lambda_j\bigr)^2=1$, with equality iff there is a single non-zero eigenvalue, we get the \emph{purity test}:
\[
  \Tr\rho^2\le1,\quad\text{with equality if and only if }\rho\text{ is pure.}
\]

\paragraph{The Bloch ball.}
Any $2\times2$ Hermitian matrix of trace $1$ is $\rho=\tfrac12(\Id+\vec r\cdot\vec\sigma)$ with $\vec r\in\R^3$ (Theorem~\ref{thm:paulibasis}), and its eigenvalues are $\tfrac12(1\pm\abs{\vec r})$. Positivity therefore means $\abs{\vec r}\le1$: \emph{qubit density operators are exactly the points of the unit ball},
with pure states on the boundary sphere (Proposition~\ref{prop:bloch}) and the maximally mixed state $\tfrac12\Id$ at the centre (Figure~\ref{fig:blochball}). Mixing states averages their Bloch vectors, and $\Tr\rho^2=\tfrac12(1+\abs{\vec r}^2)$.

\begin{figure}[htb]
  \centering
  \begin{tikzpicture}[>=stealth,scale=1.3]
    \draw[ink,thick,fill=navy!5] (0,0) circle (2);
    \draw[ink!50,dashed] (-2,0) arc (180:0:2 and 0.6);
    \draw[ink!60] (-2,0) arc (180:360:2 and 0.6);
    \draw[->,ink!70] (0,0) -- (0,2.5) node[above] {$z$};
    \fill[navy] (0,2) circle (2pt) node[above left,font=\small] {$\ket0$};
    \fill[navy] (0,-2) circle (2pt) node[below left,font=\small] {$\ket1$};
    \fill[ink] (0,0) circle (2pt) node[below left,font=\small] {$\tfrac12\Id$};
    \draw[very thick,figacc,->] (0,0) -- (0.9,0.9);
    \fill[figacc] (0.9,0.9) circle (2pt) node[right,font=\small] {mixed: $\abs{\vec r}<1$};
    \fill[navy] (1.28,1.5) circle (2pt) node[above right,font=\small] {pure: $\abs{\vec r}=1$};
  \end{tikzpicture}
  \caption{The Bloch ball. Pure states are on the surface; mixed states are inside, with the maximally mixed state $\Id/2$ at the centre.}
  \label{fig:blochball}
\end{figure}

\begin{example}\label{ex:mixture-bloch}
  Find the density matrix and Bloch vector of the ensemble $\bigl\{\tfrac12,\ket0;\ \tfrac12,\ket+\bigr\}$. Is it pure?
\end{example}
\begin{solution}
  \begin{steps}
    \item \emph{The matrix.} $\rho=\tfrac12\ket0\!\bra0+\tfrac12\ket+\!\bra+=\tfrac12\begin{pmatrix}1&0\\0&0\end{pmatrix}+\tfrac14\begin{pmatrix}1&1\\1&1\end{pmatrix}=\tfrac14\begin{pmatrix}3&1\\1&1\end{pmatrix}$.
    \item \emph{The Bloch vector.} Averaging $(0,0,1)$ and $(1,0,0)$ gives $\vec r=(\tfrac12,0,\tfrac12)$, with $\abs{\vec r}=\tfrac1{\sqrt2}<1$: mixed.
    \item \emph{Eigenvalues.} $\tfrac12(1\pm\tfrac1{\sqrt2})\approx0.85,\ 0.15$.
    \item \emph{Purity check.} $\Tr\rho^2=\tfrac1{16}(9+1+1+1)=\tfrac34=\tfrac12(1+\tfrac12)$.
  \end{steps}
\end{solution}

\begin{remark}
  A given $\rho$ arises from many different ensembles, as Example~\ref{ex:mixed-vs-superposition} showed. It can be shown that two ensembles $\{p_i,\ket{\psi_i}\}$ and $\{q_j,\ket{\varphi_j}\}$ give the same $\rho$ exactly when $\sqrt{p_i}\ket{\psi_i}=\sum_ju_{ij}\sqrt{q_j}\ket{\varphi_j}$ for a unitary matrix $(u_{ij})$ (padding the shorter ensemble with zero vectors). This is the \emph{unitary freedom} of ensembles.
\end{remark}

\subsection{States as functionals: the Riesz theorem for operators}

There is a second, more abstract, way to see what a state is, and it uses the Riesz theorem of Section~\ref{sec:state-evolution} on a space different from $\Hil$. What an experimentalist ultimately has is a rule that assigns to every observable $A$ its expectation value $\omega(A)$. That rule is a \emph{linear functional on the space of operators}, and it is positive and normalised:
\begin{steps}
  \item \emph{linear:} $\omega(aA+bB)=a\,\omega(A)+b\,\omega(B)$;
  \item \emph{positive:} $\omega(A)\ge0$ whenever $A\ge0$ (an event cannot have negative expected count);
  \item \emph{normalised:} $\omega(\Id)=1$.
\end{steps}

\begin{theorem}[Riesz theorem for operators]\label{thm:riesz-operators}
  Let $\omega:L(\Hil)\to\C$ be a linear functional. Then there is a unique operator $\rho$ such that $\omega(A)=\Tr(\rho A)$ for all $A$. Moreover, $\omega$ is positive and normalised if and only if $\rho$ is a density operator.
\end{theorem}
\begin{proof}
  \begin{steps}
    \item \emph{A Hilbert space of operators.} $L(\Hil)$ with the Hilbert--Schmidt inner product $\langle T,A\rangle=\Tr(T^\dagger A)$ is a finite-dimensional Hilbert space (Exercise~\ref{ex:hs-inner}).
    \item \emph{Apply Riesz.} By Theorem~\ref{thm:riesz} there is a unique $T$ with $\omega(A)=\langle T,A\rangle=\Tr(T^\dagger A)$. Put $\rho=T^\dagger$. This proves existence and uniqueness.
    \item \emph{Normalisation.} $\omega(\Id)=\Tr\rho$, so $\omega(\Id)=1$ iff $\Tr\rho=1$.
    \item \emph{Positivity of $\rho$.} If $\omega$ is positive, then for every $\ket\psi$ the operator $\ket\psi\!\bra\psi$ is positive, so $0\le\omega(\ket\psi\!\bra\psi)=\Tr(\rho\ket\psi\!\bra\psi)=\bra\psi\rho\ket\psi$. Thus $\rho\ge0$.
    \item \emph{Conversely,} if $\rho\ge0$ and $A\ge0$, then $\Tr(\rho A)=\Tr\bigl(\sqrt\rho A\sqrt\rho\bigr)\ge0$ because $\sqrt\rho A\sqrt\rho\ge0$ (Exercise~\ref{ex:conjugation-positive}) and positive operators have non-negative trace.
    \item \emph{Conclusion.} By Theorem~\ref{thm:density-char}, $\omega$ is positive and normalised iff $\rho$ is a density operator. \qedhere
  \end{steps}
\end{proof}

So ``a state'' has three equivalent descriptions: a unit vector up to phase (pure states), a density operator, and a positive normalised functional $A\mapsto\Tr(\rho A)$ on observables. The last is the starting point of the algebraic approach to quantum theory, in which states are \emph{defined} as such functionals.

\begin{example}\label{ex:riesz-state}
  Find the density operator representing the functional $\omega(A)=\tfrac12\bra0A\ket0+\tfrac12\bra+A\ket+$ and check that it is a state.
\end{example}
\begin{solution}
  \begin{steps}
    \item \emph{Rewrite with the trace.} $\bra\psi A\ket\psi=\Tr(\ket\psi\!\bra\psi A)$, so $\omega(A)=\Tr\bigl(\tfrac12(\ket0\!\bra0+\ket+\!\bra+)A\bigr)$.
    \item \emph{The representer.} $\rho=\tfrac12\ket0\!\bra0+\tfrac12\ket+\!\bra+=\tfrac14\begin{pmatrix}3&1\\1&1\end{pmatrix}$, the state of Example~\ref{ex:mixture-bloch}.
    \item \emph{Check.} $\Tr\rho=1$ and $\rho\ge0$ (eigenvalues $0.85,0.15$), so $\omega$ is positive and normalised.
  \end{steps}
\end{solution}

\begin{lookahead}
  The converse question, ``which probability assignments to measurement outcomes are consistent?'', has a remarkable answer, Gleason's theorem: in dimension $\ge3$, any assignment of probabilities to projectors that is additive on orthogonal projectors must be of the form $P\mapsto\Tr(\rho P)$. Probabilities \emph{are} density operators, and the Born rule is not an extra assumption.
\end{lookahead}

\subsection{The reduced density operator}

The deepest use of $\rho$ is describing a \emph{part} of a composite system. Let $\rho^{AB}$ be a density operator on $\Hil_A\tp\Hil_B$.

\begin{definition}[Partial trace and reduced density operator]\label{def:ptrace}
  The \term{partial trace} over $B$ is the linear map determined by
  \[
    \Tr_B\bigl(\ket{a_1}\!\bra{a_2}\tp\ket{b_1}\!\bra{b_2}\bigr)=\ket{a_1}\!\bra{a_2}\;\braket{b_2}{b_1}.
  \]
  The \term{reduced density operator} of $A$ is $\rho^A=\Tr_B\rho^{AB}$.
\end{definition}

It is the right description of $A$ alone because it reproduces every statistic of measurements on $A$.

\begin{proposition}\label{prop:reduced}
  For every observable $M$ on $A$, $\ \Tr\bigl((M\tp\Id)\rho^{AB}\bigr)=\Tr\bigl(M\rho^A\bigr)$.
\end{proposition}
\begin{proof}
  \begin{steps}
    \item \emph{Reduce to simple tensors.} Both sides are linear in $\rho^{AB}$, so it is enough to check on $\rho^{AB}=\ket{a_1}\!\bra{a_2}\tp\ket{b_1}\!\bra{b_2}$.
    \item \emph{Left side.} $\Tr\bigl(M\ket{a_1}\!\bra{a_2}\bigr)\Tr\bigl(\ket{b_1}\!\bra{b_2}\bigr)=\bra{a_2}M\ket{a_1}\braket{b_2}{b_1}$.
    \item \emph{Right side.} $\Tr\bigl(M\ket{a_1}\!\bra{a_2}\bigr)\braket{b_2}{b_1}$, the same. \qedhere
  \end{steps}
\end{proof}

For a pure state $\ket\psi=\sum_{ij}c_{ij}\ket{ij}$ with coefficient matrix $C=(c_{ij})$, a direct computation gives the handy formula $\rho^A=CC^\dagger$ (and $\rho^B=(C^\dagger C)^t$).

\begin{example}\label{ex:reduced-bell}
  Find the reduced density operators of each qubit of every Bell state, and of $\ket a\ket b$.
\end{example}
\begin{solution}
  \begin{steps}
    \item \emph{The state.} For $\ket{\Phi^+}=\bell00$, $\rho^{AB}=\tfrac12\bigl(\ket{00}\!\bra{00}+\ket{00}\!\bra{11}+\ket{11}\!\bra{00}+\ket{11}\!\bra{11}\bigr)$.
    \item \emph{Partial trace over $B$.} The cross terms contain $\braket10=0$ or $\braket01=0$ and vanish, while $\braket00=\braket11=1$: $\rho^A=\tfrac12\bigl(\ket0\!\bra0+\ket1\!\bra1\bigr)=\tfrac12\Id$. Likewise $\rho^B=\tfrac12\Id$.
    \item \emph{The other Bell states.} Their coefficient matrices $C=U/\sqrt2$ are $1/\sqrt2$ times a unitary $U$, so $CC^\dagger=\tfrac12\Id$.
    \item \emph{A product state.} For $\ket a\ket b$: $\rho^A=\ket a\!\bra a\,\braket bb=\ket a\!\bra a$, pure.
  \end{steps}
\end{solution}

This is a remarkable fact. The two qubits together are in a \emph{pure} state, with no uncertainty whatsoever, yet each qubit by itself is maximally mixed, as random as it can be. All the information is in the correlations between the parts. The whole is more definite than its parts, and \emph{the state of a subsystem of an entangled system is mixed}.

\begin{example}\label{ex:reduced-3term}
  Compute $\rho^A$ for $\ket\psi=\frac1{\sqrt3}(\ket{00}+\ket{01}+\ket{10})$ and its eigenvalues.
\end{example}
\begin{solution}
  \begin{steps}
    \item \emph{Coefficient matrix.} $C=\frac1{\sqrt3}\begin{pmatrix}1&1\\1&0\end{pmatrix}$.
    \item \emph{Reduced state.} $\rho^A=CC^\dagger=\tfrac13\begin{pmatrix}2&1\\1&1\end{pmatrix}$, with $\Tr\rho^A=1$.
    \item \emph{Eigenvalues.} The characteristic polynomial is $\lambda^2-\lambda+\tfrac19$, so $\lambda=\tfrac12\bigl(1\pm\tfrac{\sqrt5}3\bigr)=\tfrac{3\pm\sqrt5}6\approx0.873,\ 0.127$.
    \item \emph{Conclusion.} Both are positive and below $1$: $A$ is mixed, so $\ket\psi$ is entangled, which agrees with Example~\ref{ex:entangled-or-not}(c). (These eigenvalues are the squares of the Schmidt coefficients of Section~\ref{sec:schmidt}, and, as we will see in Section~\ref{sec:vonneumann}, determine the entanglement entropy.)
  \end{steps}
\end{solution}

\begin{example}[Reduced state of a product of mixed states]\label{ex:reduced-product}
  Let $\rho^{AB}=\rho_1\otimes\rho_2$ with $\Tr\rho_1=\Tr\rho_2=1$. Find $\rho^A$ and $\rho^B$.
\end{example}
\begin{solution}
  \begin{steps}
    \item \emph{Use linearity.} Write $\rho_1=\sum_{a,a'}r_{aa'}\ket a\!\bra{a'}$ and $\rho_2=\sum_{b,b'}s_{bb'}\ket b\!\bra{b'}$, so $\rho_1\otimes\rho_2=\sum r_{aa'}s_{bb'}\ket a\!\bra{a'}\otimes\ket b\!\bra{b'}$.
    \item \emph{Trace out $B$.} Each term contributes $r_{aa'}s_{bb'}\ket a\!\bra{a'}\braket{b'}{b}=r_{aa'}s_{bb}\delta_{bb'}\ket a\!\bra{a'}$, and summing over $b$ gives $\Tr\rho_2\cdot\rho_1=\rho_1$.
    \item \emph{Result.} $\rho^A=\rho_1$ and, symmetrically, $\rho^B=\rho_2$: the parts of a product keep their own states.
  \end{steps}
\end{solution}

\begin{exerciseset}
  \exercise{1} Compute the density matrices of the ensembles $\bigl\{\tfrac12,\ket0;\tfrac12,\ket1\bigr\}$, $\bigl\{\tfrac12,\ket+;\tfrac12,\ket-\bigr\}$ and $\bigl\{\tfrac12,\ket{+\ii};\tfrac12,\ket{-\ii}\bigr\}$. Is there a measurement that tells them apart?

  \exercise{1} A source emits $\ket0$ with probability $\tfrac23$ and $\ket{-}$ with probability $\tfrac13$. Find $\rho$, its eigenvalues, and the probability of the outcome $+1$ for $Z$ and for $X$.

  \exercise{2} Show that the density operator of an ensemble transforms as $\rho\mapsto U\rho U^\dagger$ and, under a measurement, $\rho\mapsto M_m\rho M_m^\dagger/p(m)$ with $p(m)=\Tr(M_m^\dagger M_m\rho)$ (Proposition~\ref{prop:postulates-rho}).

  \exercise{2}\label{ex:positive-hermitian-sec} Show that a positive operator is Hermitian. (Hint: write $A=B+\ii C$ with $B,C$ Hermitian and use $\bra\psi A\ket\psi\in\R$.)

  \exercise{2} Prove the purity test: $\Tr\rho^2\le1$ with equality iff $\rho$ is pure. Show that for a qubit $\Tr\rho^2=\tfrac12(1+\abs{\vec r}^2)$.

  \exercise{2} Find the density matrix, Bloch vector and purity of $\bigl\{\tfrac13,\ket0;\ \tfrac23,\ket1\bigr\}$ and of $\bigl\{\tfrac12,\ket0;\ \tfrac12,\ket+\bigr\}$ (compare Example~\ref{ex:mixture-bloch}).

  \exercise{3} Show that a state $\rho=\tfrac12(\Id+\vec r\cdot\vec\sigma)$ evolves under $U=e^{-\ii\theta\vec n\cdot\vec\sigma/2}$ to the state whose Bloch vector is $\vec r$ rotated by $\theta$ about $\vec n$. (Hint: it suffices to check for $\vec n=\hat z$.)

  \exercise{2}\label{ex:hs-inner} Show that $\langle A,B\rangle=\Tr(A^\dagger B)$ is an inner product on $L(\Hil)$, and that $\{\tfrac1{\sqrt2}\sigma_\mu\}_{\mu=0}^3$ is an orthonormal basis of $L(\C^2)$. Use it to compute the representer, in the sense of Theorem~\ref{thm:riesz-operators}, of $\omega(A)=\tfrac12\Tr(A)$.

  \exercise{2}\label{ex:conjugation-positive} Show that if $A\ge0$ then $CAC^\dagger\ge0$ for every operator $C$, and that a positive operator has non-negative trace. Deduce that $\Tr(\rho A)\ge0$ whenever $\rho,A\ge0$.

  \exercise{2} Find the density operator representing the functional $\omega(A)=\tfrac13\bra0A\ket0+\tfrac23\bra1A\ket1$, and the one representing $\omega(A)=\bra{+\ii}A\ket{+\ii}$. Which are pure?

  \exercise{3} Show that a linear functional $\omega$ with $\omega(\Id)=1$ that is \emph{not} positive can give a ``negative expectation value''. Construct one: $\omega(A)=\Tr(\rho A)$ with $\rho=\diag(\tfrac32,-\tfrac12)$, and find a positive $A$ with $\omega(A)<0$.

  \exercise{2} Let $\ket\psi=\ket a\ket b$ be a product state. Show $\rho^A=\ket a\!\bra a$. Compute the reduced density operators of each qubit for $\ket{+}\ket{0}$, the four Bell states, and $\frac1{\sqrt2}(\ket{00}+\ket{01})$.

  \exercise{2} Prove $\rho^A=CC^\dagger$ for $\ket\psi=\sum c_{ij}\ket{ij}$, and show $\rho^B=(C^\dagger C)^t$. Deduce that $\rho^A$ is pure if and only if $\rank C=1$, that is, iff $\ket\psi$ is a product state.

  \exercise{3} Show that if $\ket{\psi}$ is a pure state of $AB$ then $\rho^A$ and $\rho^B$ have the same non-zero eigenvalues. (We will see why in Section~\ref{sec:schmidt}.)

  \exercise{2} Let $\rho^{AB}=\rho_1\otimes\rho_2$. Show that $\rho^A=\rho_1$ (Example~\ref{ex:reduced-product}) and compute $\rho^A,\rho^B$ for $\rho^{AB}=\tfrac12\ket{00}\!\bra{00}+\tfrac12\ket{11}\!\bra{11}$. Is this state a product?

  \exercise{3} Let Charlie prepare $\bell00$ and send one qubit to Alice and one to Bob. Alice applies any unitary $U$ to hers. Show that Bob's reduced density operator is unchanged, so that Alice cannot send Bob a message just by acting on her qubit.
\end{exerciseset}
